\scanpage{89}
% archive sheet number: 88
l'équation $\alpha^2+\alpha-1=0$ devient $-1=0$, qui n'est jamais vérifiée (sur une base non vide !). Cela \deleted{signifie}\inserted{implique} que le nb \unsure{de|des} \deleted{\unsure{pts est}}\inserted{sommets est} au moins six, donc égal à six ou douze. Si c'était \deleted{douze} six, le stabilisateur de \deleted{\unsure{$E$}}$s_0$ serait un groupe d'ordre $20$ contenant le groupe diédral \inserted{$D_5$} -- or il n'y en a qu'un seul, qui est le produit de $D_5$ par le centre $\unsure{\mathfrak{z}|z}=\{1,\underline{a}\}$\inserted{antipodisme}. On voit que cela signifierait que $as_0=s_0$, ou $as_3=s_3$ (cela revient au $=$ par homogénéité).

Nous allons donc calculer les antipodiques des $6$ points déjà marqués (\unsure{pour faire rapide}) -- ce qui nous donnera une liste des $12$ pts !

On va d'abord déterminer le \uline{centre} de l'icosaèdre, défini (s'il existe) comme l'unique (?) pt fixe \unsure{ss|sous} $G$.
